\section{Conclusion and Future Work}

In this paper, we showed that SAEs trained on an activation that uses relative magnitude-based composition reliably learn graded special latents whose activation magnitudes correlate with $\Delta\alpha$, and demonstrated that steering these latents swaps the model's predicted output for over 50\% of inputs. 
These results indicate that the activation values of SAE latents, and not only their active/inactive status, carry information that is necessary for understanding model behaviour in this setting, which is a different perspective on transformer features than is currently prevalent in the mechanistic interpretability literature.
We release our trained SAEs on Weights \& Biases to support future work and reproducibility of our results.


\textbf{Limitations.}
The base transfomer is trained on a single minimal toy model with a narrow, controlled task. The architecture omits components standard in pretrained LLMs (MLPs, LayerNorm, multi-head attention, value and output projections) and uses a vocabulary of only 100 symbols. However, our SAE experiments demonstrate robust findings across diverse configurations: special latents (characterized by $\Delta\alpha$ correlation) emerge across three distinct SAE architectures (BatchTopK, JumpReLU, Matryoshka) with varying sparsity and dictionary size settings \ref{tab:sae_sweep_configs}. This suggests the mechanism may be fundamental to how SAEs represent the feature composition in this model, though whether relative magnitude-based composition appears in less constrained models remains an open question.

Additionally, the steering pipeline currently handles one special latent at a time; simultaneous multi-latent steering could cause even more outputs to be swapped, as suggested by the co-activation experiments in Section~\ref{ss:steering_exp}. Finally, our SAE checkpoint was selected on the basis of having few dead latents, which introduces selection bias toward high-quality SAEs. While generalisation experiments (Section~\ref{ss:sae_generalisation}) show key findings hold broadly, a fully representative study would require analysis across random samples of all trained SAEs.

\textbf{Future Work.}
The most important next step is determining whether pretrained LLMs implement similar mechanisms. The GemmaScope suites of SAEs \citep{lieberum2024gemma, mcdougall2025gemmascope2} provide a natural test environment, as it includes JumpReLU and Matryoshka SAEs trained on Gemma 2 and Gemma 3 at multiple widths; the same SAE variants we study here. The specific target would be identifying SAE latents in pretrained models where activation magnitude, rather than feature presence, is causally relevant for downstream behaviour.

Moreover, the steering pipeline could be extended in two directions: improving computational efficiency to allow exhaustive testing across all trained SAE configurations rather than a cherry-picked sample, and testing whether scaling multiple latents simultaneously can recover the remaining inputs for which single-latent steering fails. 
Preliminary evidence in Section~\ref{ss:steering_exp} suggests this may be possible. 

Finally, \cite{farrell2025order} did not fully characterise the failure modes of the base transformer model (roughly 9\% of inputs), and it is worth investigating whether these relate systematically to the cases where latent steering also fails.

In summary, the findings in this paper are intended as a foundation for further research on magnitude-based feature composition mechanisms in LLMs.